\chapter{MATLAB Dasar}

\CP{Mahasiswa mampu menggunakan MATLAB untuk membuat variabel, percabangan, pengulangan, fungsi, dan grafik sederhana serta memahami ukuran galat sebelum menerapkannya pada algoritma numerik.}

\Target{Satu skrip MATLAB dasar yang memuat variabel, \texttt{if--else}, \texttt{for}/\texttt{while}, pemanggilan fungsi, tabel hasil, dan grafik sederhana.}

\section{Variabel dan operasi dasar}

MATLAB dapat digunakan sebagai kalkulator numerik sekaligus bahasa pemrograman.
Contoh paling dasar:

\begin{lstlisting}[language=Matlab]
clear; clc; close all

a = 1;
b = 1.5;
tol = 1e-6;
maxit = 100;

x = 0:0.1:2;
y = x.^3 - x - 1;
\end{lstlisting}

Untuk operasi elemen-per-elemen pada vektor gunakan
\texttt{.*}, \texttt{./}, dan \texttt{.\^}.

\section{Percabangan \texttt{if--else}}

Percabangan digunakan ketika program harus mengambil keputusan.

\begin{lstlisting}[language=Matlab]
f = @(x) x.^3 - x - 1;

if f(a)*f(b) < 0
    disp('Interval mengurung akar')
else
    disp('Belum ada perubahan tanda')
end
\end{lstlisting}

\section{Pengulangan}

\subsection*{\texttt{for}}

Gunakan \texttt{for} ketika jumlah pengulangan sudah diketahui.

\begin{lstlisting}[language=Matlab]
for k = 1:5
    fprintf('Iterasi ke-%d\n',k)
end
\end{lstlisting}

\subsection*{\texttt{while}}

Gunakan \texttt{while} ketika pengulangan bergantung pada suatu kondisi.

\begin{lstlisting}[language=Matlab]
err = 1;
k = 0;

while err > 1e-6 && k < 100
    k = k + 1;
    err = err/2;
end
\end{lstlisting}

\section{Membuat fungsi}

Fungsi dibuat agar algoritma dapat dipakai kembali.

\begin{lstlisting}[language=Matlab]
% File: nilaiF.m
function y = nilaiF(x)
    y = x.^3 - x - 1;
end
\end{lstlisting}

Kemudian dapat dipanggil dengan

\begin{lstlisting}[language=Matlab]
fx = nilaiF(1.25);
\end{lstlisting}

\section{Menggambar grafik sederhana}

Grafik berguna untuk melihat bentuk fungsi sebelum melakukan perhitungan numerik.

\begin{lstlisting}[language=Matlab]
x = linspace(0.5,1.7,300);
y = x.^3 - x - 1;

plot(x,y,'LineWidth',1.5)
yline(0,'--')
grid on
xlabel('x')
ylabel('f(x)')
title('Grafik f(x) = x^3 - x - 1')
\end{lstlisting}

Untuk grafik konvergensi, gunakan \texttt{semilogy} jika galat turun beberapa orde besaran.

\begin{lstlisting}[language=Matlab]
errs = [1e-1 4e-2 1.2e-2 3e-3 8e-4];
semilogy(1:numel(errs),errs,'o-')
grid on
xlabel('Iterasi')
ylabel('Galat')
title('Grafik konvergensi')
\end{lstlisting}

\section{Galat: definisi sebelum latihan}

Metode numerik menghasilkan hampiran. Karena itu ukuran galat harus didefinisikan sebelum digunakan.

Jika nilai sebenarnya $x^\ast$ diketahui, galat absolut adalah

\begin{equation}
E_t = |x^\ast-x_a|.
\end{equation}

Galat relatif sebenarnya adalah

\begin{equation}
\varepsilon_t =
\frac{|x^\ast-x_a|}{|x^\ast|},
\qquad x^\ast\neq0.
\end{equation}

Pada proses iterasi, nilai sebenarnya biasanya belum diketahui.
Gunakan perubahan relatif berskala

\begin{equation}
e_a^{(k)}
=
\frac{|x_{\mathrm{new}}-x_{\mathrm{old}}|}
{\max\left(1,|x_{\mathrm{new}}|\right)}.
\end{equation}

Bentuk ini tetap terdefinisi ketika $x_{\mathrm{new}}=0$ atau sangat dekat nol.
Jika ingin dinyatakan sebagai persen, gunakan $100e_a^{(k)}\%$.

Untuk vektor,

\begin{equation}
e_a^{(k)}
=
\max_i
\frac{|x_{i,\mathrm{new}}-x_{i,\mathrm{old}}|}
{\max\left(1,|x_{i,\mathrm{new}}|\right)}.
\end{equation}

Implementasi MATLAB:

\begin{lstlisting}[language=Matlab]
function e = relerr(xnew,xold)
    e = abs(xnew-xold)./max(1,abs(xnew));
end
\end{lstlisting}

\section{Latihan MATLAB Dasar}

\begin{enumerate}
\item Buat variabel \texttt{a}, \texttt{b}, \texttt{tol}, dan \texttt{maxit}, lalu tampilkan dengan \texttt{fprintf}.
\item Buat vektor $x=0:0.1:2$ dan hitung $f(x)=x^3-x-1$.
\item Gunakan \texttt{if--else} untuk memeriksa apakah $f(a)f(b)<0$.
\item Gunakan \texttt{for} untuk menghitung kuadrat bilangan 1 sampai 10.
\item Gunakan \texttt{while} yang berhenti jika galat lebih kecil dari $10^{-6}$.
\item Buat fungsi untuk $f(x)=x^3-x-1$.
\item Buat fungsi \texttt{relerr.m} berdasarkan formula perubahan relatif berskala di atas.
\item Gambar grafik $f(x)=x^3-x-1$ dan tambahkan garis $y=0$.
\item Buat grafik \texttt{semilogy} untuk suatu vektor \texttt{errs}.
\end{enumerate}

\section{Latihan fondasi dari Lembar Kerja}

\begin{enumerate}
\item[\textbf{B0.1 (L1)}] Jelaskan apa yang dimaksud dengan \textit{metode numerik}. Apa perbedaan utama antara solusi eksak/analitik dan solusi numerik?
\item[\textbf{B0.2 (L2)}] Nilai sebenarnya adalah \(x^\ast=\sqrt{3}\) dan hampiran yang diperoleh \(x_a=1.73\). Hitung galat absolut dan galat relatif sebenarnya.
\item[\textbf{B0.3 (L3)}] Suatu iterasi menghasilkan urutan hampiran \(2.0000,\ 1.7500,\ 1.7143,\ 1.7321\). Hitung perubahan relatif antariternasi dan bandingkan formula biasa dengan formula berskala yang menggunakan \(\max(1,|x_{\mathrm{new}}|)\).
\item[\textbf{B0.4 (L4)}] Jelaskan mengapa dua hampiran dapat mempunyai galat absolut yang sama tetapi kualitasnya berbeda. Kaitkan dengan galat relatif.
\item[\textbf{B0.5 (L5)}] Jelaskan peran \textit{round-off error}, \textit{truncation error}, toleransi, dan jumlah iterasi maksimum dalam algoritma numerik.
\end{enumerate}
